\section{Reverse tile-shape inference}
\label{sec:inference}

\subsection{The question, and the two halves of the answer}

Welder's \textsc{Propagate} starts from a chosen \emph{output} tile and walks the
graph backwards. At each node it asks which region of each input the tile needs.
Two separate mechanisms answer this:

\begin{enumerate}
\item \textbf{Within a node:} \code{resolve_inputs(spec, output_tile, consts)} in
  \code{kernels/teeny-kernels/src/graph/tile_resolve.rs}. This is a pure function
  that is fully implemented and tested.
\item \textbf{Across an edge:} the input tile a consumer needs becomes the output
  tile its producer must supply. This is the graph walk. It is not implemented
  today (\code{TileGraph} has only \code{from_dag}). \S\ref{sec:edges} describes what
  it must do and where the DSL makes that hard.
\end{enumerate}

A tile is a \code{TileShape = Vec<Option<usize>>}: one entry per \emph{real} dim,
with \code{None} meaning ``unresolved, therefore full extent''.

\subsection{\texttt{resolve\_inputs}, step by step}

\begin{lstlisting}
fn resolve_inputs(spec, output_tile, consts) -> Result<Vec<TileShape>> {
    let out = spec.outputs[0];                     // error if there are no outputs
    require output_tile.len() == out.rank;          // error otherwise -- the BN2d trap, Sec. 7
    // 1. read the blocks the output was given, keyed by NAME
    let resolved = for axis in out.axes:
        block = product of output_tile[d] for d in axis.dims   // None if any is None
        if block is Some: (axis.extent_param, block)
    // 2. for each input, start all-None and place only what can be placed
    for input in spec.inputs:
        tile = [None; input.rank]
        for axis in input.axes:
            name  = axis.window.map(|w| w.output_extent_param).unwrap_or(axis.extent_param)
            block = match resolved.get(name) {
                Some(b) => b,
                None if axis.window.is_some() => continue,         // stays full extent
                None => consts.get(axis.extent_param) / divide_by, // or continue
            };
            extent = match axis.window {
                None => block,
                Some(w) => (block - 1) * stride(w) + kernel(w),    // receptive field; or continue
            };
            tile[last of axis.dims] = extent;  tile[other dims] = 1;  // product-preserving
        yield tile
}
\end{lstlisting}

The rules, each of which has a test in \code{tile_resolve.rs}:

\begin{description}[leftmargin=1.2em]
\item[Names, not arithmetic.] An input axis takes the block of the output axis with
  the same \code{extent_param}. GEMM's \code{a}/\code{c} both name \code{M}, and
  transpose's permuted axes resolve by name regardless of position. An axis whose
  name the output never declares, such as GEMM's \code{K}, stays \code{None}. That
  is correct, because a reduction axis is the search's to choose, not the output's.
\item[Unresolved means full extent.] This is the conservative direction: an
  over-estimate is safe against a hard memory capacity, and an under-estimate is
  not. It covers three cases:
  \begin{itemize}
  \item an axis with no binding (untiled, walked or reduced);
  \item an unmatched name;
  \item a window whose consts the caller did not supply.
  \end{itemize}
\item[Windows invert by the receptive field.] A windowed input axis resolves against
  the \emph{output} variable its window names (\code{OW}), not its own extent
  (\code{W}). The tile is then $(b-1)\cdot s + k$. Padding shifts the origin, not
  the size, so it never enters. The pad family is the degenerate window
  $s=k=1$, whose tile is $b$.
\item[A window never takes the full-extent fallback through the window.] Feeding
  \code{H} through $(H-1)s+k$ would claim the kernel reads past the tensor. An
  unresolved windowed axis therefore stays \code{None}.
\item[Only the output's \emph{bound} axes seed anything.] The output tile has a value
  in every dim, but \code{resolve_inputs} reads only the dims covered by
  \code{out.axes}. A value at an untiled output dim (\code{B}, \code{C_OUT},
  \code{OH} for conv2d) is ignored. This single rule produces conv2d's large
  over-estimate (\S\ref{sec:conv-overestimate}).
\item[Flattened bindings are product-preserving.] The block lands on the innermost
  dim of \code{axis.dims} and the others become \code{1}. Reading a flattened
  \emph{output} binding multiplies its dims back together, so the two directions are
  inverse \emph{on the product}. They are not an axis-aligned region once the block
  exceeds the innermost extent (\S\ref{sec:edges}).
\item[\texttt{divide\_by} replaces a full extent.] It is applied only on the
  \code{consts} fallback path, never to a propagated block. The macro always emits
  \code{None}, because GroupNorm's \code{C / G} is per-instance
  (\issue{1nr.15}).
\end{description}

\subsection{Crossing an edge}
\label{sec:edges}

The resolved input tile is in the consumer's \emph{dim coordinates}. The producer's
\code{resolve_inputs} needs an output tile in \emph{its} dim coordinates. Names do
not help across an edge: conv2d's output calls its axes \code{B/C_OUT/OH/OW},
BatchNorm2d's input \code{B/C/HW}, and SiLU's \code{n_elements}. The only shared
coordinate system is the tensor's real dims, so the walk must pass \code{TileShape}s
\emph{by position}. That imposes three requirements on the DSL.

\begin{enumerate}
\item \textbf{Ranks must agree.} A kernel whose view of a tensor has a different rank
  from the edge cannot take part. That was BatchNorm2d before \code{span}.
\item \textbf{Full extent must be materialised.} \code{None} on the consumer side
  means ``all of it''. The producer must read that as the real extent, which comes
  from the node's shape, not as ``unconstrained''.
\item \textbf{Flattened tiles must be re-expressed axis-aligned, or the producer will
  misread them.} SiLU's flat spec turns an output tile into \code{[1, 1, 1, 1024]}.
  BatchNorm2d's spanned \code{HW} passes that on as \code{[_, _, 1, 1024]}. Handed to
  conv2d as an output tile, that says ``1024 columns of \code{OW}'' on a tensor 32
  wide. The honest translation is \code{[_, _, 32, 32]} (32 rows of 32). Welder's own
  implementation clamps every propagated region to the producer's real shape
  (\code{graph.py}, \code{propogate_inputs}), and a teenygrad walk needs at least
  that clamp. It also needs the reshape from a product to an axis-aligned box
  whenever the block exceeds the innermost extent.
\end{enumerate}

\begin{note}[A useful invariant]
A flattened block $b$ over dims $(\ldots, H, W)$ is an axis-aligned
$\lceil b/W\rceil \times \min(b, W)$ box exactly when $b \le W$ and $W \bmod b = 0$,
or when $W \mid b$. For BatchNorm2d's \code{BLOCK_HW = 128} on a $32\times32$ plane,
that is a $4\times32$ box. A walk can convert under this condition and fall back to
full extent otherwise.
\end{note}
